% Fragmento dos Resultados (RQ07) para o template SBC (Overleaf).
% Issue #40 -- colar dentro de \section{Resultados}.
% NUMEROS PARCIAIS: gerados por `python -m analise.rq07` sobre os repositorios com coleta
% completa no momento da escrita. Atualizar tabelas e numeros do texto ao final da coleta.
%
% Bibliografia: adicione as entradas abaixo ao .bib do grupo (ou ajuste as citacoes).
% @article{cohen1968weighted, author={Cohen, Jacob}, title={Weighted kappa: nominal scale
%   agreement with provision for scaled disagreement or partial credit},
%   journal={Psychological Bulletin}, volume={70}, number={4}, pages={213--220}, year={1968}}
% @article{landis1977measurement, author={Landis, J. Richard and Koch, Gary G.},
%   title={The measurement of observer agreement for categorical data},
%   journal={Biometrics}, volume={33}, number={1}, pages={159--174}, year={1977}}

\subsection{RQ07 --- A classifica\c{c}\~ao DORA depende da defini\c{c}\~ao operacional?}

\textbf{M\'etodo.} Classificamos cada reposit\'orio em \textit{Elite}, \textit{High},
\textit{Medium} ou \textit{Low} com a tabela de refer\^encia da disciplina: cada m\'etrica
recebe uma nota de 1 a 4, e a classifica\c{c}\~ao geral \'e a mediana das quatro notas,
arredondada para baixo. Repetimos a classifica\c{c}\~ao sob tr\^es combina\c{c}\~oes de
defini\c{c}\~oes operacionais (Tabela~\ref{tab:rq07-combinacoes}). Todas usam o mesmo tempo
de recupera\c{c}\~ao, que depende apenas dos \textit{workflow runs} e n\~ao da defini\c{c}\~ao
de \textit{release}. \textbf{C1} \'e a combina\c{c}\~ao de refer\^encia. \textbf{C3} altera uma
\'unica defini\c{c}\~ao, a do \textit{lead time}, de (a) para (b). \textbf{C2} altera tr\^es de
uma vez: inclui pr\'e-\textit{releases} entre os \textit{deploys} e adota o \textit{lead time}~(b)
e o CFR~(b). Comparamos as combina\c{c}\~oes duas a duas, entre os reposit\'orios
classific\'aveis em ambas, pela propor\c{c}\~ao que muda de categoria e pelo kappa ponderado
linear~\cite{cohen1968weighted}, interpretado pela conven\c{c}\~ao de Landis e
Koch~\cite{landis1977measurement}.

\begin{table}[ht]
\centering
\caption{Combina\c{c}\~oes de defini\c{c}\~oes comparadas.}
\label{tab:rq07-combinacoes}
\small
\begin{tabular}{lllll}
\hline
Combina\c{c}\~ao & Unidade de \textit{deploy} & \textit{Lead time} & CFR & Recupera\c{c}\~ao \\
\hline
C1 (refer\^encia) & \textit{releases} est\'aveis & (a) & (a) & mediana dos epis\'odios \\
C2 & \textit{releases} e pr\'e-\textit{releases} & (b) & (b) & mediana dos epis\'odios \\
C3 & \textit{releases} est\'aveis & (b) & (a) & mediana dos epis\'odios \\
\hline
\end{tabular}
\end{table}

\textbf{Resultados.} A an\'alise cobre 164 reposit\'orios com coleta completa; 149 deles
(91\%) t\^em as quatro m\'etricas e s\~ao classific\'aveis nas tr\^es combina\c{c}\~oes. A
Tabela~\ref{tab:rq07-dist} mostra a distribui\c{c}\~ao das categorias. Sob C1, 49,0\% dos
reposit\'orios s\~ao \textit{Medium} e 40,9\% s\~ao \textit{High} ou \textit{Elite}. Sob C2,
nenhum reposit\'orio \'e \textit{Elite} e o n\'umero de reposit\'orios \textit{Low} passa de 15
para 25. C3 fica pr\'oxima de C1.

\begin{table}[ht]
\centering
\caption{Distribui\c{c}\~ao das categorias por combina\c{c}\~ao ($n = 149$).}
\label{tab:rq07-dist}
\small
\begin{tabular}{lrrrr}
\hline
Combina\c{c}\~ao & \textit{Elite} & \textit{High} & \textit{Medium} & \textit{Low} \\
\hline
C1 & 10 (6,7\%) & 51 (34,2\%) & 73 (49,0\%) & 15 (10,1\%) \\
C2 & 0 (0,0\%) & 35 (23,5\%) & 89 (59,7\%) & 25 (16,8\%) \\
C3 & 7 (4,7\%) & 57 (38,3\%) & 74 (49,7\%) & 11 (7,4\%) \\
\hline
\end{tabular}
\end{table}

A Tabela~\ref{tab:rq07-pares} compara as combina\c{c}\~oes. Mudar apenas a defini\c{c}\~ao de
\textit{lead time} (C1 e C3) altera a categoria de 18,1\% dos reposit\'orios, com
concord\^ancia substancial ($\kappa = 0{,}751$). Combinar as tr\^es mudan\c{c}as (C1 e C2) altera
a categoria de 45,0\% dos reposit\'orios, com concord\^ancia apenas razo\'avel
($\kappa = 0{,}341$). C2 e C3 diferem em 37,6\% dos casos ($\kappa = 0{,}464$, concord\^ancia
moderada).

\begin{table}[ht]
\centering
\caption{Concord\^ancia entre combina\c{c}\~oes ($n = 149$).}
\label{tab:rq07-pares}
\small
\begin{tabular}{lrrl}
\hline
Par & \% que mudou de categoria & $\kappa$ ponderado & Interpreta\c{c}\~ao \\
\hline
C1 e C2 & 45,0\% & 0,341 & razo\'avel \\
C1 e C3 & 18,1\% & 0,751 & substancial \\
C2 e C3 & 37,6\% & 0,464 & moderada \\
\hline
\end{tabular}
\end{table}

\textbf{Discuss\~ao.} A classifica\c{c}\~ao DORA de um reposit\'orio n\~ao \'e neutra em rela\c{c}\~ao
\`a defini\c{c}\~ao operacional adotada. Trocar s\'o a forma de agregar o \textit{lead time}
preserva a categoria de 82\% dos reposit\'orios, enquanto a combina\c{c}\~ao de unidade de
\textit{deploy}, \textit{lead time} e CFR diferentes reclassifica quase metade deles e elimina a
categoria \textit{Elite} nesta amostra. Isso indica que compara\c{c}\~oes entre estudos, ou com os
relat\'orios de refer\^encia, s\'o s\~ao interpret\'aveis quando a defini\c{c}\~ao operacional \'e
declarada, e que uma an\'alise de sensibilidade deveria acompanhar qualquer classifica\c{c}\~ao
baseada em dados de reposit\'orios.

\textbf{Limita\c{c}\~oes.} (i) \textit{Cobertura parcial.} Os n\'umeros se referem aos 164
reposit\'orios com coleta completa na data da escrita, todos nas quatro faixas de estrelas mais
baixas (1.000 a 25.000 estrelas; 42, 42, 41 e 39 reposit\'orios por faixa). As tr\^es faixas
superiores e os 12 reposit\'orios com mais \textit{releases} ainda n\~ao estavam inclu\'idos; os
resultados n\~ao devem ser extrapolados para a amostra completa. (ii) \textit{Estratifica\c{c}\~ao.}
A amostra tem aloca\c{c}\~ao igual por faixa de estrelas, ent\~ao as propor\c{c}\~oes descrevem a
amostra analisada e n\~ao a popula\c{c}\~ao de reposit\'orios do GitHub. (iii) \textit{Desenho n\~ao
fatorial.} Como C2 altera tr\^es defini\c{c}\~oes simultaneamente, n\~ao \'e poss\'ivel atribuir a
mudan\c{c}a de 45,0\% a uma delas em particular; s\'o a defini\c{c}\~ao de \textit{lead time} \'e
isolada (C1 e C3). (iv) \textit{Desvio em rela\c{c}\~ao \`a metodologia.} A variante por
\textit{tags}, prevista como unidade de \textit{deploy}, n\~ao foi coletada por custo de chamadas
\`a API; C3 a substitui variando apenas o \textit{lead time}. (v) \textit{Validade do CFR~(b).} A
heur\'istica de \textit{release} corretiva ainda n\~ao foi validada contra a amostra-ouro, de modo
que C2 herda essa incerteza. (vi) \textit{Reposit\'orios n\~ao classific\'aveis.} 15 dos 164
reposit\'orios n\~ao t\^em todas as m\'etricas e ficam de fora: 14 n\~ao t\^em epis\'odio de falha
conclu\'ido no CI (nenhuma falha, ou falhas sem recupera\c{c}\~ao at\'e o fim da janela), o que torna o
tempo de recupera\c{c}\~ao ausente por defini\c{c}\~ao, e um n\~ao tem \textit{lead time} calcul\'avel. Como esses
reposit\'orios saem da classifica\c{c}\~ao, os resultados valem apenas para os que falharam e se
recuperaram ao menos uma vez, o que pode subestimar as categorias mais altas.
